\section{The final dihedral endpoint}
\label{sec:r2053-endpoint}

\begin{proposition}\label{prop:r2053-r18-reduction}
A complex categorification of $R7\_4\_2053$ would produce a complex
categorification of $R8\_1\_18$.
\end{proposition}
\begin{proof}
Use Proposition~\ref{prop:four-dihedral}. The source has integral
subcategory of Tambara--Yamagami type $C_4$, and its two irrational
simples are self-dual. Its rank-sixteen Morita dual has
\[
 D_8=\langle r,s\mid r^4=s^2=1,\ srs=r^{-1}\rangle,
 \qquad \rho r=r^{-1}\rho,\quad \rho s=rs\rho,
 \qquad \rho^2=a+N\rho.
\]
Here $a\in\{1,r^2\}$ and $N$ is the sum of the rotations or the sum
of the reflections. In the inverse reflection dual, the two irrational
supports are the size-four double orbits
$\langle s\rangle g\phi(\langle s\rangle)$, with trivial stabilizers.
Duality preserves each orbit when $a=1$ and interchanges them when
$a=r^2$. Since each support therefore carries exactly one simple
bimodule, source self-duality forces $a=1$.

Suppose first that $N=P=1+r+r^2+r^3$. The two-dimension theorem
\cite[Lemma~3.1]{Schopieray} permits a pseudo-unitary Galois conjugate
of the actual Morita context. In its canonical positive spherical
structure, $\rho$ is self-dual and
\[
 \rho^3=P+(1+4P)\rho.
\]
For a self-dual simple $W$, a form of sign $\nu_2(W)$ induces a form
of the same sign on $W^{\otimes3}$: the dual reverses three identical
factors. Each self-dual constituent occurring oddly has the same
indicator, since a nondegenerate alternating multiplicity-space form
has even dimension. The unit and $r^2$ both occur once above, hence
$\nu_2(\rho)=\nu_2(r^2)=1$.

The pointed associator on $\langle r\rangle=C_4$ consequently has
even class in $H^3(C_4,\C^\times)\simeq\Z/4$, because restriction
to $\langle r^2\rangle$ is reduction modulo two. Its restriction
therefore has order at most two. The integral inverse reflection dual
is group-theoretical. Its Frobenius--Schur exponent, computed as the
least common multiple of
$|C|\operatorname{ord}(\omega|_C)$ over cyclic subgroups $C\leq D_8$,
divides eight: the rotation term divides eight and every reflection
term divides four. But this category is of type $\operatorname{TY}(C_4)$.
For either nondegenerate bicharacter
$\chi(x,y)=i^{\varepsilon xy}$, $\varepsilon=\pm1$, the eighth
indicator of its dimension-two simple is
\[
 \nu_8(X)=\frac18
 \sum_{a+b+c+d=0\ ({\rm mod}\ 4)}
 i^{\varepsilon(ab+ac+ad+bc+bd+cd)}=0.
\]
Indeed the four exponent classes occur $8,24,8,24$ times. This holds
for either sign of the Tambara--Yamagami parameter. Exponent dividing
eight would instead give $\nu_8(X)=\FPdim(X)=2$, a contradiction.
The formulas used here are \cite[Theorem~9.2]{NgSchauenburgExponent}
and \cite[Theorem~3.4]{ShimizuTY}.

Thus $N=s+rs+r^2s+r^3s$. Put $\sigma=s\rho$. The mixed rules give
\[
 \sigma r=r\sigma,\qquad
 \sigma^2=r^3+(1+r+r^2+r^3)\sigma.
\]
The full additive subcategory on
$1,r,r^2,r^3,\sigma,r\sigma,r^2\sigma,r^3\sigma$ is closed under
products and duals. Reversing the rotation generator identifies its
fusion ring with $R8\_1\_18$.
\end{proof}

The independent source verifier
\texttt{verify\_r2053\_r18\_reduction.py} reopens both census tensors,
checks all source identities and both complete rank-sixteen fusion
tensors, recomputes the exact indicator sums, and matches all $512$
coefficients of the closed rank-eight subcategory. Its local-to-census
map is $(0,1,7,2,4,3,6,5)$. Thus the endpoint is an actual fusion
subcategory of the Morita dual, not merely a dimension match.


\begin{theorem}\label{thm:r2053-final}
The fusion ring $R7\_4\_2053$ is not the Grothendieck ring of a
complex fusion category.
\end{theorem}
\begin{proof}
Proposition~\ref{prop:r2053-r18-reduction} reduces the question to
$R8\_1\_18$. The general exclusion of that ring is proved in the
companion rank-eight classification \cite{Rank8}; its complete
source-bound proof is
\texttt{R18\_STS\_MOMENT\_PROOF.md} in the rank-eight evidence.
We recall the mechanism to make the dependency explicit.

For that ring, the two-dimension theorem gives a pseudo-unitary
Galois conjugate. Write $t=2+\sqrt5$ and $D=8+16t$.
Cauchy's theorem restricts the Frobenius--Schur exponent to the
primes $2,5$. Positivity of induction for the pointed $C_4$
subcategory gives a four-point Fourier decomposition and bounds
its generator's multiplicity in each central forgetful row by one.
Character reciprocity leaves eight complete induced-unit frames.
For the normalized center modular data, let $C$ be the source
character matrix and $V$ the corresponding induced-unit forgetful
rows. The twist-one identity $STS=T^{-1}ST^{-1}$ yields
\[
 \sum_A v_Av_A^{\mathsf T}\theta_A
       =V^{\mathsf T}\overline C.
\]
Relative trace to $\mathbb Q(\sqrt5)$ gives a fixed first-moment
Gram target. First and second induction moments, together with
rationality of the relative fourth indicators, give the remaining
constraints. The functional $a+bt\mapsto2a-b$ annihilates $D$
times any rational fourth indicator. Prime support makes the
relative moment vectors from orders dividing forty, together with
the zero vector, exhaustive; there are thirteen in all. Full Galois row
orbits and the four possible pointed associator classes are retained.

Three integer covectors, one for each effective pointed class
(the two odd classes have identical constraints), are nonnegative
on every admissible column and have pairing $-256$ with every target.
The exact source verifier checks all eight frames in each class;
its $24$ branch contradictions use no numerical solver.
Thus $R8\_1\_18$ cannot be categorified. The actual fusion
subcategory supplied by Proposition~\ref{prop:r2053-r18-reduction}
would contradict this theorem.
\end{proof}

Together with the preceding exclusions and explicit models,
Theorem~\ref{thm:r2053-final} completes the rank-seven census through
multiplicity four: all $4201$ source rings are classified, with
$4172$ general-complex exclusions and $29$ unitary realizations.
